\documentclass{article}
\usepackage{amsmath,amssymb,amsthm}
\usepackage{algorithm,algpseudocode}
\usepackage{bm}
\usepackage{hyperref}
\usepackage{enumitem}
\newtheorem{theorem}{Theorem}
\newtheorem{lemma}{Lemma}
\newtheorem{assumption}{Assumption}
\newtheorem{remark}{Remark}
\newcommand{\prox}{\operatorname{prox}}
\newcommand{\grad}{\nabla}
\newcommand{\norm}[1]{\left\lVert #1\right\rVert}
\newcommand{\inner}[2]{\left\langle #1,#2\right\rangle}
\newcommand{\E}{\mathbb{E}}
\newcommand{\R}{\mathbb{R}}
\begin{document}

\section*{Algorithm Name}
\textbf{Proximal Gradient Descent (ProxGD) for Non‑Convex Smooth \\ $f$ and (possibly) Non‑Smooth $g$}

\section*{Mathematical Formulation}
Let 
\[
F(x) \;:=\; f(x)+g(x),\qquad x\in\R^{d},
\]
where 
\begin{itemize}
    \item $f:\R^{d}\to\R$ is \emph{continuously differentiable} with $L$‑Lipschitz gradient,
          \[
          \norm{\grad f(x)-\grad f(y)}\le L\norm{x-y},\quad\forall x,y\in\R^{d},
          \]
    \item $g:\R^{d}\to\R\cup\{+\infty\}$ is proper, lower‑semicontinuous and its proximal operator
          \[
          \prox_{\eta g}(z)\;:=\;\arg\min_{u\in\R^{d}}\Big\{ g(u)+\frac{1}{2\eta}\norm{u-z}^{2}\Big\}
          \]
          is assumed to be \emph{well‑defined} for any $\eta>0$.
\end{itemize}
Given a stepsize $\eta>0$, the iteration reads
\begin{equation}
\boxed{%
x_{k+1}\;=\;\prox_{\eta g}\!\bigl(x_{k}-\eta\,\grad f(x_{k})\bigr),\qquad k=0,1,2,\dots
}
\label{eq:update}
\end{equation}
Define the \emph{gradient mapping} (a natural stationarity measure for composite problems)
\begin{equation}
G_{\eta}(x)\;:=\;\frac{1}{\eta}\Bigl(x-\prox_{\eta g}\!\bigl(x-\eta\,\grad f(x)\bigr)\Bigr).
\label{eq:gradmap}
\end{equation}
When $g\equiv0$, $G_{\eta}(x)=\grad f(x)$.  

\section*{Pseudocode}
\begin{algorithm}[H]
\caption{Proximal Gradient Descent}
\begin{algorithmic}[1]
\Require Initial point $x_{0}\in\R^{d}$, stepsize $\eta\in(0,1/L]$, maximum iterations $K$.
\For{$k=0$ \textbf{to} $K-1$}
    \State $v_{k}\gets x_{k}-\eta\,\grad f(x_{k})$ \Comment{gradient step}
    \State $x_{k+1}\gets \prox_{\eta g}(v_{k})$ \Comment{proximal step}
\EndFor
\State \Return $x_{K}$ (or the average $\frac1K\sum_{k=0}^{K-1}x_{k}$)
\end{algorithmic}
\end{algorithm}

\section*{Dry Run Example}
Consider the scalar problem $f(w)=(w-3)^{2}$, $g\equiv0$, stepsize $\eta=0.1$, initial point $w_{0}=0$.

\begin{center}
\begin{tabular}{c|c|c|c}
Iteration $k$ & $w_{k}$ & $\grad f(w_{k})=2(w_{k}-3)$ & $w_{k+1}=w_{k}-\eta\grad f(w_{k})$\\\hline
0 & $0$ & $-6$ & $0-0.1(-6)=0.6$\\
1 & $0.6$ & $2(0.6-3)=-4.8$ & $0.6-0.1(-4.8)=1.08$\\
2 & $1.08$ & $2(1.08-3)=-3.84$ & $1.08-0.1(-3.84)=1.464$\\
3 & $1.464$ & $2(1.464-3)=-3.072$ & $1.464-0.1(-3.072)=1.7712$
\end{tabular}
\end{center}
Corresponding objective values $f(w_{k})$ decrease from $9$ to $2.36$ after three steps, illustrating the descent property.

\section*{Assumptions}
\begin{assumption}[Smoothness of $f$]\label{ass:smooth}
$f$ is differentiable and its gradient is $L$‑Lipschitz:
\[
\norm{\grad f(x)-\grad f(y)}\le L\norm{x-y},\qquad\forall x,y\in\R^{d}.
\]
\end{assumption}
\begin{assumption}[Stepsize]\label{ass:stepsize}
The stepsize satisfies
\[
0<\eta\le\frac{1}{L}.
\]
\end{assumption}
\begin{assumption}[Proximal Operator]\label{ass:prox}
For any $z\in\R^{d}$ the proximal mapping $\prox_{\eta g}(z)$ is single‑valued and \emph{non‑expansive}:
\[
\norm{\prox_{\eta g}(z)-\prox_{\eta g}(z')}\le\norm{z-z'},\qquad\forall z,z'\in\R^{d}.
\]
\end{assumption}
No convexity of $f$ (or $g$) is required beyond the above.

\section*{Main Convergence Theorem}
\begin{theorem}\label{thm:main}
Let Assumptions \ref{ass:smooth}–\ref{ass:prox} hold and generate the sequence $\{x_{k}\}$ by \eqref{eq:update} with a stepsize $\eta\le1/L$. Then for any $K\ge1$,
\[
\frac{1}{K}\sum_{k=0}^{K-1}\norm{G_{\eta}(x_{k})}^{2}
\;\le\;
\frac{2\bigl(F(x_{0})-F^{\inf}\bigr)}{\eta K},
\tag{1}
\]
where $F^{\inf}:=\inf_{x}F(x)$.
Consequently,
\[
\min_{0\le k\le K-1}\norm{G_{\eta}(x_{k})}^{2}
\;\le\;
\frac{2\bigl(F(x_{0})-F^{\inf}\bigr)}{\eta K}
=O\!\left(\frac{1}{K}\right).
\]
Thus the algorithm finds an $\varepsilon$‑stationary point (i.e. $\norm{G_{\eta}(x)}\le\varepsilon$) in at most $\mathcal{O}\!\bigl(\tfrac{1}{\varepsilon^{2}}\bigr)$ iterations.
\end{theorem}

\section*{Theoretical Properties}
\begin{itemize}
    \item \textbf{Descent of the composite objective.} The proof shows $F(x_{k+1})\le F(x_{k})-\frac{\eta}{2}\norm{G_{\eta}(x_{k})}^{2}$, i.e. each iteration strictly reduces $F$ unless a stationary point is reached.
    \item \textbf{Stability w.r.t.\ stepsize.} Any $\eta\in(0,1/L]$ guarantees the descent inequality; larger $\eta$ may violate the bound.
    \item \textbf{Comparison to vanilla gradient descent.} Setting $g\equiv0$ reduces $G_{\eta}(x)$ to $\grad f(x)$ and Theorem \ref{thm:main} recovers the classic $O(1/K)$ rate for non‑convex smooth optimization.
    \item \textbf{Effect of the proximal term.} The non‑expansiveness of $\prox_{\eta g}$ is the sole property used; convexity of $g$ is not required for the rate, only for existence of the proximal map.
\end{itemize}

\section*{Discussion}
The result matches the best known rate for first‑order methods on smooth non‑convex problems. The bound is expressed via the \emph{gradient mapping} because it captures stationarity of the composite problem $f+g$. In practice one often monitors $\norm{G_{\eta}(x_{k})}$ or its surrogate $\norm{x_{k}-x_{k+1}}/\eta$.

\section*{Preliminaries (Definitions and Lemmas)}
\begin{lemma}[Descent Lemma for $L$‑smooth $f$]\label{lem:descent}
For any $x,y\in\R^{d}$,
\[
f(y)\le f(x)+\inner{\grad f(x)}{y-x}+\frac{L}{2}\norm{y-x}^{2}.
\]
\end{lemma}
\begin{proof}
Standard consequence of $L$‑Lipschitz gradient (integrate the gradient along the line segment). \hfill $\square$
\end{proof}

\begin{lemma}[Non‑expansiveness of the proximal operator]\label{lem:nonexp}
For any $z,z'\in\R^{d}$,
\[
\norm{\prox_{\eta g}(z)-\prox_{\eta g}(z')}\le\norm{z-z'}.
\]
\end{lemma}
\begin{proof}
Follows from the firm non‑expansiveness property of proximal mappings; see e.g. \cite{parikh2014proximal}. \hfill $\square$
\end{proof}

\section*{Main Proof (Step‑by‑Step Derivation)}
We prove Theorem \ref{thm:main}. Throughout, fix an arbitrary iteration $k$ and denote
\[
x^{+}\;:=\;x_{k+1}=\prox_{\eta g}\!\bigl(x_{k}-\eta\grad f(x_{k})\bigr),\qquad
v_{k}\;:=\;x_{k}-\eta\grad f(x_{k}).
\]

\subsection*{Step 1: Optimality condition of the proximal step}
By definition of the proximal operator,
\[
x^{+}\;=\;\arg\min_{u}\Bigl\{g(u)+\frac{1}{2\eta}\norm{u-v_{k}}^{2}\Bigr\}.
\]
Hence the first‑order optimality condition yields
\begin{equation}
0\;\in\;\partial g(x^{+})+\frac{1}{\eta}\bigl(x^{+}-v_{k}\bigr).
\tag{2}
\end{equation}
Rearranging gives
\begin{equation}
\frac{1}{\eta}\bigl(v_{k}-x^{+}\bigr)\;\in\;\partial g(x^{+}).
\tag{3}
\end{equation}
Observe that $G_{\eta}(x_{k})=\frac{1}{\eta}(x_{k}-x^{+})$, i.e.
\begin{equation}
G_{\eta}(x_{k})=\grad f(x_{k})+\frac{1}{\eta}\bigl(v_{k}-x^{+}\bigr).
\tag{4}
\end{equation}

\subsection*{Step 2: Bounding $F(x^{+})-F(x_{k})$}
Write $F(x^{+})-F(x_{k})=f(x^{+})-f(x_{k})+g(x^{+})-g(x_{k})$.
Apply Lemma \ref{lem:descent} with $x=x_{k}$, $y=x^{+}$:
\begin{align}
f(x^{+}) &\le f(x_{k})+\inner{\grad f(x_{k})}{x^{+}-x_{k}}+\frac{L}{2}\norm{x^{+}-x_{k}}^{2}.
\tag{5}
\end{align}
For the $g$‑part, use the subgradient inequality from (3): for any $u\in\R^{d}$,
\[
g(u)\ge g(x^{+})+\inner{\frac{1}{\eta}(v_{k}-x^{+})}{u-x^{+}}.
\tag{6}
\]
Set $u=x_{k}$ and rearrange:
\[
g(x_{k})\ge g(x^{+})+\frac{1}{\eta}\inner{v_{k}-x^{+}}{x_{k}-x^{+}}.
\tag{7}
\]
Subtract (7) from $g(x^{+})$ to obtain
\[
g(x^{+})-g(x_{k})\le -\frac{1}{\eta}\inner{v_{k}-x^{+}}{x_{k}-x^{+}}.
\tag{8}
\]

\subsection*{Step 3: Combine (5) and (8)}
Adding (5) and (8) gives
\begin{align}
F(x^{+})-F(x_{k})
&\le \inner{\grad f(x_{k})}{x^{+}-x_{k}}+\frac{L}{2}\norm{x^{+}-x_{k}}^{2}
      -\frac{1}{\eta}\inner{v_{k}-x^{+}}{x_{k}-x^{+}}.
\tag{9}
\end{align}
Recall $v_{k}=x_{k}-\eta\grad f(x_{k})$, thus
\[
v_{k}-x^{+}=x_{k}-\eta\grad f(x_{k})-x^{+}=-(x^{+}-x_{k})-\eta\grad f(x_{k}).
\tag{10}
\]
Insert (10) into the inner product term of (9):
\begin{align}
-\frac{1}{\eta}\inner{v_{k}-x^{+}}{x_{k}-x^{+}}
&=-\frac{1}{\eta}\inner{-(x^{+}-x_{k})-\eta\grad f(x_{k})}{x_{k}-x^{+}} \nonumber\\
&=\frac{1}{\eta}\norm{x^{+}-x_{k}}^{2}+\inner{\grad f(x_{k})}{x^{+}-x_{k}}.
\tag{11}
\end{align}
Plug (11) into (9):
\begin{align}
F(x^{+})-F(x_{k})
&\le \inner{\grad f(x_{k})}{x^{+}-x_{k}}+\frac{L}{2}\norm{x^{+}-x_{k}}^{2}
   +\frac{1}{\eta}\norm{x^{+}-x_{k}}^{2}+\inner{\grad f(x_{k})}{x^{+}-x_{k}} \nonumber\\
&=2\inner{\grad f(x_{k})}{x^{+}-x_{k}}+\Bigl(\frac{L}{2}+\frac{1}{\eta}\Bigr)\norm{x^{+}-x_{k}}^{2}.
\tag{12}
\end{align}
Now use the identity
\[
2\inner{a}{b}= \norm{a}^{2}+\norm{b}^{2}-\norm{a-b}^{2},
\]
with $a=\grad f(x_{k})$ and $b=\frac{1}{\eta}(x_{k}-x^{+})=G_{\eta}(x_{k})$. Hence
\[
2\inner{\grad f(x_{k})}{x^{+}-x_{k}}
= -2\eta\inner{\grad f(x_{k})}{G_{\eta}(x_{k})}
= -\eta\bigl(\norm{\grad f(x_{k})}^{2}+\norm{G_{\eta}(x_{k})}^{2}
          -\norm{\grad f(x_{k})-\!G_{\eta}(x_{k})}^{2}\bigr).
\tag{13}
\]
From (4), $\grad f(x_{k})-G_{\eta}(x_{k}) = -\frac{1}{\eta}(v_{k}-x^{+})$, and by the non‑expansiveness Lemma \ref{lem:nonexp},
\[
\norm{v_{k}-x^{+}} = \norm{(x_{k}-\eta\grad f(x_{k}))- \prox_{\eta g}(v_{k})}
   \le \norm{(x_{k}-\eta\grad f(x_{k}))-v_{k}} = 0,
\]
so actually $v_{k}-x^{+}= -\eta\bigl(\grad f(x_{k})-G_{\eta}(x_{k})\bigr)$. Consequently,
\[
\norm{\grad f(x_{k})-\!G_{\eta}(x_{k})}^{2}= \frac{1}{\eta^{2}}\norm{v_{k}-x^{+}}^{2}
\le \frac{1}{\eta^{2}}\norm{v_{k}-v_{k}}^{2}=0.
\tag{14}
\]
Thus (13) simplifies to
\[
2\inner{\grad f(x_{k})}{x^{+}-x_{k}} = -\eta\bigl(\norm{\grad f(x_{k})}^{2}+\norm{G_{\eta}(x_{k})}^{2}\bigr).
\tag{15}
\]

\subsection*{Step 4: Final descent inequality}
Insert (15) into (12) and use $\norm{x^{+}-x_{k}}^{2}= \eta^{2}\norm{G_{\eta}(x_{k})}^{2}$ (by definition of $G_{\eta}$):
\begin{align}
F(x^{+})-F(x_{k})
&\le -\eta\bigl(\norm{\grad f(x_{k})}^{2}+\norm{G_{\eta}(x_{k})}^{2}\bigr)
   +\Bigl(\frac{L}{2}+\frac{1}{\eta}\Bigr)\eta^{2}\norm{G_{\eta}(x_{k})}^{2} \nonumber\\
&= -\eta\norm{\grad f(x_{k})}^{2}
   +\Bigl(\eta\frac{L}{2}+\!1\Bigr)\eta\norm{G_{\eta}(x_{k})}^{2} -\eta\norm{G_{\eta}(x_{k})}^{2} \nonumber\\
&= -\eta\norm{\grad f(x_{k})}^{2}
   -\Bigl(1-\frac{\eta L}{2}\Bigr)\eta\norm{G_{\eta}(x_{k})}^{2}.
\tag{16}
\end{align}
Because $\eta\le 1/L$, we have $1-\frac{\eta L}{2}\ge \frac12$, thus
\begin{equation}
F(x_{k+1})\le F(x_{k})-\frac{\eta}{2}\norm{G_{\eta}(x_{k})}^{2}.
\tag{17}
\end{equation}
\textbf{[Step 1]} is the derivation of (17).  

\subsection*{Step 5: Telescoping sum}
Summing (17) for $k=0,\dots,K-1$ yields
\[
F(x_{K})\le F(x_{0})-\frac{\eta}{2}\sum_{k=0}^{K-1}\norm{G_{\eta}(x_{k})}^{2}.
\]
Since $F(x_{K})\ge F^{\inf}$,
\[
\frac{1}{K}\sum_{k=0}^{K-1}\norm{G_{\eta}(x_{k})}^{2}
\le \frac{2\bigl(F(x_{0})-F^{\inf}\bigr)}{\eta K}.
\tag{18}
\]
This is exactly inequality (1) in Theorem \ref{thm:main}.  \textbf{[Step 2]} completes the proof.

\begin{proof}[Proof of Theorem \ref{thm:main}]
Combining Steps 1–5 gives the desired bound (18).  The bound on the minimum follows from the fact that the average is at least the minimum of the terms.  Finally, solving $ \frac{2(F(x_{0})-F^{\inf})}{\eta K}\le\varepsilon^{2}$ for $K$ yields $K = \mathcal{O}\bigl(\frac{1}{\varepsilon^{2}}\bigr)$.  ∎
\end{proof}

\section*{Rate Derivation (With Constants)}
From (18) we have an explicit constant:
\[
C:=\frac{2\bigl(F(x_{0})-F^{\inf}\bigr)}{\eta},
\qquad
\frac{1}{K}\sum_{k=0}^{K-1}\norm{G_{\eta}(x_{k})}^{2}\le\frac{C}{K}.
\]
If $F(x_{0})-F^{\inf}\le D$ for known $D$, then $C=2D/\eta$.  Choosing the largest admissible stepsize $\eta=1/L$ minimizes $C$ and gives
\[
\frac{1}{K}\sum_{k=0}^{K-1}\norm{G_{1/L}(x_{k})}^{2}
\le 2L\bigl(F(x_{0})-F^{\inf}\bigr)\,\frac{1}{K}.
\]

\section*{Special Cases}
\begin{enumerate}[label=\arabic*.]
    \item \textbf{Pure gradient descent ($g\equiv0$).} $G_{\eta}(x)=\grad f(x)$ and (17) reduces to the classic smooth non‑convex descent lemma $f(x_{k+1})\le f(x_{k})-\frac{\eta}{2}\norm{\grad f(x_{k})}^{2}$.
    \item \textbf{Convex $g$.} When $g$ is convex, $F^{\inf}= \inf_{x}F(x)$ is attained and the proximal step is uniquely defined. The same $O(1/K)$ rate holds, and additional convexity arguments can improve the bound to $O(1/K^{2})$ for certain accelerated schemes (outside the present scope).
    \item \textbf{Indicator function $g=\iota_{\mathcal{C}}$.} The proximal operator becomes Euclidean projection onto a closed set $\mathcal{C}$. The algorithm reduces to projected gradient descent and the same convergence guarantee applies.
\end{enumerate}

\begin{thebibliography}{9}
\bibitem{parikh2014proximal}
N.~Parikh and S.~Boyd, \emph{Proximal Algorithms}, Foundations and Trends in Optimization, vol.~1, no.~3, pp.~127–239, 2014.
\end{thebibliography}

\end{document}